\section{Discussion and Limitations}
\label{sec:discussion}

\paragraph{Why use the nonlinear law?}
The proposed surface provides input-dependent increment shaping,
rather than a uniform increase in calibrated gains. The strict
reference comparison demonstrates that this change can reduce
transient buffer demand while preserving certified convergence and
execution constraints. Matched-gain Callisto startup tests support a
similar measured effect under local delayed execution. Independently
tuned tests answer a different practical question: the selected
Fuzzy scheme, including its chosen gains, improves disagreement RMS
in the two declared random regimes. Keeping these comparisons
separate avoids attributing every benefit to inference alone.

\paragraph{Engineering interpretation.}
Lower excursion increases reserve at fixed capacity; it does not
demonstrate a hardware-capacity saving. Since $\gamma_i=100/(d_iC)$,
changing capacity or degree changes the controller map unless its
normalization/gains are adjusted. Similarly, lower disagreement RMS
can accompany a larger common-frequency shift, because there is no
absolute-rate reference loop. Command TV measures activity, not
energy. The application must identify acceptable buffer reserves,
frequency shift and correction bounds when choosing between tasks.

\paragraph{Scope of the evidence.}
The certificate is restricted to its synchronous zero-latency fluid
assumptions and invariant domain. Callisto uses a continuous
plant with asynchronous local updates and quantized frame events;
finite observations do not prove its infinite-horizon safety.
Uncertified does not mean unsafe, and saturation does not by itself
mean buffer failure. Persistent random inputs are assessed through
finite-window deviations rather than exact nominal equilibrium.
Frequency agreement, positive progress and capacity compliance
support normal model execution, but do not replace an application
refinement proof of functional logical synchrony.

An important distinction is between a controller's regional increment
bound and a closed-loop storage guarantee. The rule-average property
$|F|\le5$ limits one increment, whereas the accumulator can retain
past action. The reference certificate closes the input domain and
then bounds disagreement, common frequency and actuator use jointly.
The experimentally observed nonlinear memory is consistent with this
stateful behavior, but its same-input reconstruction is not a separate
closed-loop ablation. Thus the analysis supplies a reason to investigate
the response, while the independent tests supply the evidence for its
usefulness in the tested regimes.

The safety sweep also separates proof conservatism from execution
constraints. Failure of a sufficient inequality at $s=4$ leaves the
trajectory unclassified by that certificate; the finite observations
remain useful without relabeling them as an all-time guarantee.
At $s=32$, limiter intervention requires a different interpretation
from the non-clipped theorem, even though measured buffers remain
within capacity. Neither the number of feasible grid candidates nor
one favorable initial-frequency direction establishes robust safe-set
inclusion over arbitrary disturbances or initial states.

The empirical scope is a three-node triangle, a fixed gain grid and
Fuzzy shape, and one training seed per regime. Independent tests
quantify variation within the specified input families, not global
optimality or arbitrary-network robustness. Ordered batch transport
omits finite serialization, dropout and hardware execution costs.
Observers and FFP switching are disabled. Extending the invariant
analysis to local sampling, bounded delay and quantization, followed
by broader-topology and hardware validation, is the next step.
Appendix~\ref{app:discussion} gives the fuller interpretation and
the requirements for further controller extensions.
